\section{Cơ sở lý thuyết}
\label{sec:background}

\subsection{Leo thang đặc quyền và thoát container trong Kubernetes}
Trong cụm Kubernetes, các container/pod được cô lập mức logic bằng namespace,
cgroup và capabilities, nhưng chia sẻ chung nhân của nút. Bề mặt nguy hiểm nhất là
\emph{leo thang đặc quyền} và \emph{thoát container}: tận dụng lỗ hổng nhân hoặc
cấu hình sai (capabilities dư thừa, mount host-path, \texttt{ptrace},
\texttt{unshare}/\texttt{setns}, nạp mô-đun nhân) để vượt ranh giới cô lập và chiếm
quyền trên nút~\cite{ref5_combe_docker}. Đặc điểm chung: mọi con đường khai thác đều
biểu hiện thành các \emph{chuỗi syscall đặc quyền} đặc trưng (Bảng~\ref{tab:bg-cve}),
nên phân tích syscall là điểm quan sát tự nhiên để phát hiện chúng.

\begin{table}[htbp]
  \centering
  \caption{Một số họ lỗ hổng leo thang đặc quyền/thoát container và dấu vết syscall.}
  \label{tab:bg-cve}
  \small
  \begin{tabularx}{\linewidth}{@{}l >{\raggedright\arraybackslash}p{3.6cm} >{\raggedright\arraybackslash}X@{}}
    \toprule
    \textbf{Lỗ hổng} & \textbf{Cơ chế} & \textbf{Syscall liên quan} \\
    \midrule
    PwnKit (CVE-2021-4034)   & \texttt{pkexec} setuid + tải \texttt{.so} & \texttt{execve}, \texttt{openat}/\texttt{mmap} \\
    Dirty Pipe (CVE-2022-0847) & ghi đè trang pipe    & \texttt{splice} (kèm \texttt{write}) \\
    runc escape (CVE-2019-5736) & ghi đè \texttt{/proc/self/exe} & \texttt{openat}, \texttt{execve} \\
    Lạm dụng capability        & nạp mô-đun/đặc quyền  & \texttt{init\_module}, \texttt{capset}, \texttt{mount} \\
    \bottomrule
  \end{tabularx}
\end{table}

\subsection{Giám sát lời gọi hệ thống cho container}
Các cơ chế truyền thống (seccomp, AppArmor, SELinux) lọc theo \emph{chính sách
tĩnh} nên: (1) không thích nghi với hành vi thời gian chạy; (2) thiếu nhận thức về
\emph{chuỗi} syscall nên mù với tấn công nhiều bước; và (3) tốn chi phí cấu hình thủ
công ở quy mô lớn. eBPF cho phép nạp logic giám sát/thực thi an toàn vào nhân tại
các điểm vào syscall với chi phí thấp~\cite{ref20_bpf_security}, và là nền tảng để
thu thập luồng syscall mà một mô-đun phát hiện bất thường cần. Trọng tâm của báo cáo
là chính \emph{lớp phát hiện} (mô-đun SyscallAD), độc lập với cơ chế thu thập cụ thể;
việc thu thập syscall thời gian thực được xem là đầu vào cho sẵn.

\subsection{Phát hiện xâm nhập dựa trên chuỗi syscall}
Hướng phát hiện bất thường dựa trên chuỗi syscall có lịch sử lâu dài, khởi nguồn từ
``a sense of self''~\cite{forrest1996sense,hofmeyr1998ids} và các mô hình thay
thế~\cite{warrender1999detecting}, cùng cảnh báo kinh điển về tấn công bắt
chước~\cite{wagner2002mimicry}. Các hướng hiện đại gồm phương pháp ngữ
nghĩa~\cite{creech2014semantic}, túi lời gọi (BoSC) cho
container~\cite{abed2015bag}, và học sâu một lớp/biến phân như
DAGMM~\cite{zong2018dagmm}, Deep SVDD~\cite{ruff2018deepsvdd}, HIDS học
sâu~\cite{haider2021deep}. Riêng cho container có các đánh giá phát hiện bất
thường~\cite{castanhel2021peek}, ngữ cảnh hóa syscall~\cite{elkhairi2022chids} và
tập hợp Random/Isolation Forest~\cite{bsoul2022ensemble}. SyscallAD nằm trong dòng
này: hợp nhất VAE với iForest và bổ sung khớp mẫu lành tính bằng PrefixSpan.

\subsection{Công trình gần đây (2023--2026)}
Các nghiên cứu mới cho thấy SyscallAD nằm trong một không gian thiết kế \emph{đang
hoạt động sôi nổi}. Ở lớp thực thi, lọc syscall bằng eBPF có trạng thái đã vượt
seccomp tĩnh nhờ nhận thức ngữ cảnh và đối số~\cite{progsyscall2023}; các agent eBPF
thời gian thực như \mbox{eBPF-PATROL} phát hiện trực tiếp leo thang đặc quyền, thoát
container và reverse shell từ luồng syscall~\cite{ebpfpatrol2025}. Ở lớp phát hiện
bất thường, \emph{chính công thức} VAE\,+\,iForest của SyscallAD được tái sử dụng và
mở rộng: FedMon đưa nó lên học liên kết đa cụm Kubernetes~\cite{fedmon2025}, còn
LIGHT-HIDS ghép trích đặc trưng DeepSVDD với bộ phát hiện \emph{một lớp} bằng Rừng cô
lập~\cite{lighthids2025}; mô hình ngôn ngữ trên chuỗi syscall (LSTM/Transformer)
được dùng cho phát hiện \emph{tính mới}~\cite{lmnovelty2023}. Ngược lại, các khảo sát
gần đây~\cite{dlidssurvey2025,translmidssurvey2024} chỉ ra rằng mô hình
Transformer/LLM tuy nắm ngữ cảnh tốt hơn nhưng \emph{chi phí và độ trễ cao}, khó dùng
để chặn nội tuyến thời gian thực---củng cố lựa chọn lõi nhẹ VAE\,+\,iForest. Cuối
cùng, các phương pháp đối kháng (tiêm seqGAN, bắt chước) cho thấy bộ phát hiện theo
chuỗi vẫn có thể bị né tránh, và đề xuất hướng tăng cường độ bền
vững~\cite{advrobustseq2025}---một mối đe dọa mà chuyên đề bàn ở Mục~\ref{sec:discussion}.

\subsection{Bộ mã hóa biến phân (VAE) và Rừng cô lập (iForest)}
Hai khối học máy lõi của SyscallAD:
\begin{itemize}[leftmargin=1.6em]
  \item \textbf{VAE}~\cite{vae_kingma} học phân phối ẩn của dữ liệu \emph{bình
        thường}; điểm bất thường là \emph{lỗi tái cấu trúc}: mẫu lệch khỏi phân phối
        bình thường được tái tạo kém nên có lỗi lớn. VAE phù hợp với học một lớp
        (chỉ cần dữ liệu lành tính), thích hợp cho phát hiện tấn công zero-day.
  \item \textbf{iForest}~\cite{iforest_liu} cô lập ngoại lệ bằng phân vùng ngẫu
        nhiên đệ quy: điểm hiếm bị cô lập ở độ sâu nông hơn, cho điểm bất thường cao.
        iForest bổ trợ VAE ở khía cạnh \emph{ngoại lệ thống kê} mà lỗi tái cấu trúc
        có thể bỏ sót.
\end{itemize}

\subsection{Bộ dữ liệu DongTing}
DongTing~\cite{ref26_dongtingdataset} là bộ dữ liệu quy mô lớn gồm các chuỗi syscall
của nhân Linux được gán nhãn bình thường/độc hại (theo các bản vá lỗi syzkaller),
thích hợp để huấn luyện và kiểm chứng bộ phát hiện bất thường. Đây cũng là nguồn dữ
liệu mà SyscallAD trong DeSFAM sử dụng; báo cáo này dùng DongTing làm nguồn dữ liệu
ngoại tuyến để tái lập và đánh giá. Chuỗi bình thường lưu dưới dạng \emph{ID
syscall}; chuỗi tấn công lưu dưới dạng \emph{tên syscall}; bảng phân chia
(\texttt{Baseline.xlsx}) định nghĩa các tập \texttt{DTDS-train} /
\texttt{DTDS-validation} / \texttt{DTDS-test}.
